\documentclass[../../Main.tex]{subfiles}
\begin{document}
\section{添加多语言支持}
对于Ren'Py 8与Future分支用户来说，您可以跳过此部分内容，因为在Future分支中已内建多语言支持。

对于Ren'Py 7用户，您需要将Future分支中的部分代码移植到您的项目中，以实现多语言支持。

打开screens.rpy文件，在第一行\verb|init offset = -1|下添加以下代码：
\begin{lstlisting}
init python:
    def scan_translations():

        languages = renpy.known_languages()

        if not languages:
            return None

        rv = [(i, renpy.translate_string("{#language name and font}", i)) for i in languages ]
        rv.sort(key=lambda a : renpy.filter_text_tags(a[1], allow=[]).lower())

        rv.insert(0, (None, "简体中文"))

        bound = math.ceil(len(rv)/2.)

        return (rv[:bound], rv[bound:2*bound])

default translations = scan_translations()
\end{lstlisting}

然后，按住Ctrl + F打开搜索框，输入“选择选项之后”，另起一行，输入以下代码（注意缩进，第一行if语句应与上一段落的vbox在同一缩进级别）：
\begin{lstlisting}
                if translations:
                    vbox:
                        style_prefix "radio"
                        label _("语言")
                        hbox:
                            viewport:
                                mousewheel True
                                scrollbars "vertical"
                                ysize 100
                                has vbox

                                for tran in translations:
                                    vbox:
                                        for tlid, tlname in tran:
                                            textbutton tlname:
                                                action Language(tlid)
\end{lstlisting}

\section{生成翻译文件}
打开Ren'Py SDK，选中您的项目，在右侧操作栏中点击生成翻译，输入语言名字，例如：\verb|mylanguage|，点击生成翻译文件。稍等片刻，您便会在game/tl/目录下看到一个名为mylanguage的文件夹，其中包含了多个rpy文件，依次点开，对照old的内容将其中的文本翻译成您需要的语言，填在new里，保存。

然后，打开对应语言文件夹下的screens.rpy，在第一行\verb|translation|后另起一行添加以下代码：

\begin{lstlisting}
    old "{#language name and font}"
    new "<您的语言名称>"
\end{lstlisting}

例如：

\begin{lstlisting}
    old "{#language name and font}"
    new "My Language"
\end{lstlisting}

最后重新运行游戏，您就可以在游戏中设置菜单选择新的语言了。

\end{document}